
\chapter{Procesos Gaussianos}

\noindent
Este bloque de la teoría se centrará en \textbf{variables y vectores Gaussianos}, herramienta fundamental para el estudio del Movimiento Browniano que aparecerá como límite en la sección anterior.

\begin{definicion}[Variable Aleatoria Normal]
    Sea $(\Omega,\mathcal{F},P)$ un espacio de probabilidad. Una v.a. $X:\Omega \to \mathbb{R}$ es \textbf{normal o Gaussiana}, $X\sim N(\mu,\sigma^2)$, si admite densidad
    \[
    f_X(x) = \frac{1}{\sqrt{2\pi}\,\sigma}\, e^{-\dfrac{(x-\mu)^2}{2\sigma^2}}, \qquad \mu\in\mathbb{R},\ \sigma>0.
    \]
    Se demuestra que necesariamente $\mu=E[X]$ y $\sigma^2=\text{Var}(X)$.
\end{definicion}

\begin{definicion}[Centrado y estandarización]
    Una v.a. $Z$ es \textbf{centrada} si $E[Z]=0$. Se dice que $X\sim N(0,1)$ es \textbf{normal estándar}, donde la estandarización
    \[
    Z = \frac{X-E[X]}{\sigma} \sim N(0,1) \qquad\Longleftrightarrow\qquad X = \sigma Z + \mu \sim N(\mu,\sigma^2)
    \]
    permite reducir cualquier cálculo sobre una normal genérica a uno sobre la normal estándar.
\end{definicion}

\begin{prop}[Propiedades de las v.a. normales]
Sea $X\sim N(\mu,\sigma^2)$. Entonces:
\begin{enumerate}
    \item \textbf{Función característica:}
    \[
    \hat{f}_X(k) = E[e^{ikX}] = e^{i\mu k - \frac{1}{2}\sigma^2 k^2}.
    \]
    \item \textbf{Momentos} (caso centrado $X\sim N(0,\sigma^2)$):
    \[
    E[X^n] = \begin{cases} \dfrac{(2k)!}{2^k k!}\,\sigma^{2k}, & n=2k \\ 0, & n=2k+1\end{cases}, \qquad k=0,1,2,\dots
    \]
    es decir, todos los momentos impares se anulan por simetría y los pares se expresan en función de $\sigma$.
    \item \textbf{Concentración:} para $\kappa>0$,
    \[
    P(|X-\mu|\le \kappa\sigma) = 2\Phi(\kappa)-1,
    \]
    con $\Phi$ la f.d.a. de $Z\sim N(0,1)$. Esto cuantifica cómo se concentra la masa de probabilidad alrededor de la media en unidades de desviación típica.
    \item \textbf{Colas:} para $Z\sim N(0,1)$ y $x>0$,
    \[
    \frac{1}{\sqrt{2\pi}}\left(\frac{1}{x}-\frac{1}{x^3}\right)e^{-x^2/2} \le P(Z\ge x) \le \frac{1}{\sqrt{2\pi}}\frac{1}{x}e^{-x^2/2},
    \]
    lo que muestra el conocido decaimiento \textit{super-rápido} (gaussiano) de las colas.
\end{enumerate}
\end{prop}

\begin{teo}[Combinación lineal de normales]
Sean $X\sim N(\mu_1,\sigma_1^2)$ e $Y\sim N(\mu_2,\sigma_2^2)$ independientes, y $a\in\mathbb{R}$. Entonces
\[
X+Y \sim N(\mu_1+\mu_2,\ \sigma_1^2+\sigma_2^2), \qquad aX \sim N(a\mu_1,\ a^2\sigma_1^2).
\]
\end{teo}
\begin{observacion}[Idea de la demostración]
Por independencia, la función característica de la suma es el producto de las funciones características:
\[
\hat{f}_{X+Y}(k) = \hat{f}_X(k)\hat{f}_Y(k) = e^{i(\mu_1+\mu_2)k - \frac{1}{2}(\sigma_1^2+\sigma_2^2)k^2},
\]
que es exactamente la función característica de $N(\mu_1+\mu_2,\sigma_1^2+\sigma_2^2)$. Como la función característica determina unívocamente la ley, se concluye. El caso $aX$ es análogo evaluando $\hat{f}_X(ak)$.
\end{observacion}

\begin{prop}[Estabilidad de la normalidad bajo límites]
Sea $(X_n)_{n\ge0}$ una sucesión de v.a. normales independientes, $X_n\sim N(\mu_n,\sigma_n^2)$, que converge en distribución a una v.a. $X$. Entonces:
\begin{enumerate}
    \item El límite es también normal, $X\sim N(\mu,\sigma^2)$, con $\mu=\lim_n \mu_n$ y $\sigma^2 = \lim_n \sigma_n^2$.
    \item Si además $(X_n)$ converge a $X$ en probabilidad, entonces converge también en $L^p$ para todo $1\le p<\infty$.
\end{enumerate}
\end{prop}

\section{Vectores Aleatorios Gaussianos}

\begin{definicion}[Vector Gaussiano]
Sea $X=(X_1,\dots,X_n):\Omega\to\mathbb{R}^n$. Decimos que $X$ es un \textbf{vector Gaussiano} si \textbf{toda} combinación lineal de sus componentes es una v.a. normal:
\[
\sum_{i=1}^n c_i X_i \ \text{es normal}, \qquad \forall\, c_1,\dots,c_n\in\mathbb{R}.
\]
\end{definicion}

\begin{observacion}
Si $X$ es Gaussiano, cada componente $X_i$ es necesariamente normal (tomando $c=e_i$). \textbf{El recíproco es falso}: que todas las marginales sean normales no implica que el vector sea Gaussiano.
\end{observacion}

\begin{ejemplo}[Contraejemplo del recíproco]
Sea $X\sim N(0,1)$ y $a>0$ fijo. Definimos
\[
Y = \begin{cases} X & \text{si } |X|\le a \\ -X & \text{si } |X|> a.\end{cases}
\]
Se puede comprobar que $Y\sim N(0,1)$ (es normal), pero el par $(X,Y)$ \textbf{no} es un vector Gaussiano (por ejemplo $X+Y$ no es normal, ya que vale $2X$ en $|X|\le a$ y $0$ fuera).
\end{ejemplo}

\begin{definicion}[Media y matriz de covarianzas]
Sea $X=(X_1,\dots,X_n)$ un vector Gaussiano. Se definen:
\[
E[X] = \big(E[X_1],\dots,E[X_n]\big) \quad\text{(vector de medias)},
\]
\[
C(X) = \big(\text{Cov}(X_i,X_j)\big)_{i,j=1}^n \quad\text{(matriz de covarianzas)}, \quad \text{Cov}(X_i,X_j) = E[X_iX_j]-E[X_i]E[X_j].
\]
Un vector Gaussiano se dice \textbf{centrado} si $E[X]=0$.
\end{definicion}

\begin{teo}[Función característica de un vector Gaussiano]
    Sea $X=(X_1,\dots,X_n)$ un vector Gaussiano con media $m$ y matriz de covarianzas $C(X)$. Entonces, para todo $x\in\mathbb{R}^n$,
    \[
    \hat{f}_X(x) = \exp\left(i\langle m,x\rangle - \frac{1}{2}\langle x, C(X) x\rangle\right).
    \]
\end{teo}
\begin{observacion}[Idea de la demostración]
    Para $a=(a_1,\dots,a_n)\in\mathbb{R}^n$, por definición de vector Gaussiano la combinación lineal $\langle a,X\rangle = \sum_i a_iX_i$ es una v.a. normal, con
    \[
    E[\langle a,X\rangle] = \langle a,m\rangle, \qquad \text{Var}(\langle a,X\rangle) = \langle a, C(X) a\rangle.
    \]
    Su función característica (escalar) es entonces $e^{i\langle a,m\rangle - \frac{1}{2}\langle a,C(X)a\rangle}$. Tomando $a=x$ y evaluando en el punto $1$ se obtiene exactamente $\hat f_X(x)$.
\end{observacion}

\begin{prop}[Independencia y covarianza nula]
    Sean $X,Y$ dos v.a. normales (conjuntamente Gaussianas). Entonces
    \[
    X \perp\!\!\!\perp Y \quad \Longleftrightarrow \quad \text{Cov}(X,Y)=0.
    \]
\end{prop}
\begin{observacion}[Idea de la demostración]
    ($\Rightarrow$) Es válido en general: independencia implica $E[XY]=E[X]E[Y]$, luego $\text{Cov}(X,Y)=0$. \\
    ($\Leftarrow$) Aquí se usa que el vector es Gaussiano: si $\text{Cov}(X,Y)=0$, la matriz de covarianzas del vector $(X,Y)$ es diagonal, $C=\text{diag}(\sigma_X^2,\sigma_Y^2)$. Sustituyendo en la función característica del vector,
    \[
    \hat f_{(X,Y)}(u) = e^{-\frac{1}{2}(u_1^2\sigma_X^2+u_2^2\sigma_Y^2)} = \hat f_X(u_1)\,\hat f_Y(u_2),
    \]
    y como la función característica conjunta factoriza como producto de las marginales, $X$ e $Y$ son independientes.
\end{observacion}

\begin{teo}[Densidad Gaussiana multivariante]
    Sea $X=(X_1,\dots,X_n)$ un vector Gaussiano en $\mathbb{R}^n$ con media $m$ y matriz de covarianzas $C(X)$ \textbf{invertible}. Entonces $X$ admite densidad de probabilidad
    \[
    f_X(x) = \frac{1}{(2\pi)^{n/2}\,(\det C)^{1/2}} \; \exp\left(-\frac{1}{2}\langle x-m,\; C^{-1}(x-m)\rangle\right), \qquad x\in\mathbb{R}^n.
    \]
\end{teo}

\begin{observacion}
    Esta fórmula generaliza directamente la densidad normal univariante: cuando $n=1$ se recupera 
    $$f_X(x) = \frac{1}{\sqrt{2\pi}\sigma}e^{-(x-\mu)^2/2\sigma^2}$$
    La condición de invertibilidad de $C$ es necesaria: si $C$ es singular, el vector vive en realidad en un subespacio de dimensión menor que $n$ y no existe densidad respecto a la medida de Lebesgue en $\mathbb{R}^n$.
\end{observacion}

\section{Multivariant Gaussian Distribution}

\begin{definicion}[Función de Densidad de Probabilidad Multivariante]
Sea $\underline{X} = (X_1, \dots, X_n)$ un vector aleatorio Gaussiano en el espacio euclídeo $\mathbb{R}^n$, caracterizado por un vector de valores esperados (medias) $m = \mathbb{E}[\underline{X}]$ y una matriz de covarianzas $C(\underline{X})$. Asumiendo que la matriz de covarianzas es regular y admite inversa (es decir, existe su matriz de precisión $C^{-1}(\underline{X})$), entonces la \textbf{función de densidad de probabilidad (f.d.p.)} del vector $\underline{X}$ viene definida de forma explícita por la expresión:
\begin{equation*}
    f_{\underline{X}}(x) = \frac{1}{(2\pi)^{n/2} \left(\det(C)\right)^{1/2}} \, e^{-\frac{1}{2} \langle x-m, \, C^{-1}(x-m) \rangle}
\end{equation*}
donde $\langle \cdot, \cdot \rangle$ denota el producto escalar euclídeo estándar en $\mathbb{R}^n$, el cual parametriza la forma cuadrática de la distancia de Mahalanobis.

\begin{proof}
    % //TODO: entender demostración
    % Para construir la demostración rigurosa de la f.d.p. generalizada, estructuramos el análisis en dos etapas fundamentales: el establecimiento del caso base estándar e independiente, y la posterior aplicación de una transformación lineal geométrica mediante el uso de operadores matriciales.

    % \vspace{0.4cm}

    % \noindent \textbf{Etapa 1: El Vector Gaussiano Estándar ($m = 0$, $C = I$)} \\
    % Supongamos inicialmente que $\underline{X}$ constituye un vector Gaussiano estándar. En este escenario estacionario, las componentes son variables normales canónicas univariantes e independientes, con vector de medias nulo ($m = 0$) y una matriz de covarianzas equivalente a la matriz identidad de orden $n$:
    % \[
    % C(\underline{X}) = I_n = \begin{pmatrix} 1 & \dots & 0 \\ \vdots & \ddots & \vdots \\ 0 & \dots & 1 \end{pmatrix} \quad \text{($n$ elementos en la diagonal)}
    % \]
    % Debido a la hipótesis de independencia estocástica mutua entre las variables $X_1, \dots, X_n$, la función de densidad conjunta se factoriza exactamente como el producto ordinario de sus densidades marginales unidimensionales:
    % \begin{align*}
    %     f_{\underline{X}}(x) &= \overbrace{f_{X_1}(x_1) f_{X_2}(x_2) \dots f_{X_n}(x_n)}^{n \text{ términos}} = \frac{1}{\sqrt{2\pi}} e^{-\frac{x_1^2}{2}} \frac{1}{\sqrt{2\pi}} e^{-\frac{x_2^2}{2}} \dots \frac{1}{\sqrt{2\pi}} e^{-\frac{x_n^2}{2}} = \\
    %     &= \frac{1}{(2\pi)^{n/2}} \exp\left( -\frac{\sum_{i=1}^{n} x_i^2}{2} \right) = \frac{1}{(2\pi)^{n/2}} e^{-\frac{\|x\|^2}{2}}
    % \end{align*}
    % Donde el vector de evaluación se denota por $x = (x_1, \dots, x_n)$ y su norma euclídea al cuadrado $\|x\|^2$ se vincula directamente con el producto interno simétrico:
    % \[
    % \|x\|^2 = \sum_{i=1}^{n} x_i^2 = \langle x, x \rangle
    % \]
\end{proof}
\end{definicion}


\begin{lema}[Isomorfismo Estocástico Gausiano e Inversión Lineal]
    Si el vector aleatorio $\underline{X}$ es Gaussiano con vector de medias $m$ y una matriz de covarianzas $C(\underline{X})$ simétrica y definida positiva (y por ende, invertible), se cumple de forma unívoca que la transformación lineal normalizada:
    \[
    C^{-1/2}(\underline{X}-m) \sim \mathcal{N}(0, I_n) \quad \text{constituye un vector Gaussiano estándar (v.G.e.).}
    \]
    Inversamente, si partimos de un vector Gaussiano estándar nativo denotado como $Z \sim \mathcal{N}(0, I_n)$, la transformación afín mapeada por la raíz cuadrada matricial y el desplazamiento del promedio genera el vector original:
    \[
    \underline{X} = m + C^{1/2}(Z) \quad \text{el cual recupera las propiedades } \mathbb{E}[\underline{X}] = m \text{ y } \text{Var}(\underline{X}) = C(\underline{X}).
    \]
\end{lema}

\begin{flushright}
\textit{Representación sintética de síntesis:} $\underline{X} = C^{1/2}(Z) + m \quad \text{con } Z \sim \mathcal{N}(0, I_n)$.
\end{flushright}

\strut
\hfill \rule{\textwidth}{0.5pt} \hfill
\strut

\noindent \textbf{Etapa 2: Aplicación del Teorema del Cambio de Variable mediante Medida de Cajas} \\
Definamos un dominio canónico $E$ como cualquier hiperrectángulo o caja $n$-dimensional acotada inmersa en el espacio medible $\mathbb{R}^n$:
\[
{\color{green!50!black}E = \bigotimes_{i=1}^{n} [a_i, b_i]}
\]
Determinar que el vector de interés se localiza en el interior de dicha región equivale a restringir cada componente marginal a sus intervalos independientes correspondientes:
\[
\underline{X} \in E \equiv X_1 \in [a_1, b_1], \dots, X_n \in [a_n, b_n]
\]
Modelemos la medida de probabilidad de este suceso utilizando la relación de la transformación lineal deducida en el Lema ($\underline{X} = C^{1/2}(Z) + m$). Al operar de forma algebraica con los límites del conjunto y aplicar el formalismo clásico de integración multidimensional bajo cambios de variable coordenados, se desarrolla la siguiente cadena de identidades:
\begin{align*}
    \mathbb{P}(\underline{X} \in E) &= \mathbb{P}\left(C^{1/2}(Z) + m \in E\right) = \mathbb{P}\left[Z \in C^{-1/2}(E-m)\right] = \int_{C^{-1/2}(E-m)} f_Z(u) \, du = 
\end{align*}
Introduciendo la sustitución geométrica del cambio de base lineal $y = C^{1/2}u + m$, el diferencial de volumen en la medida de Lebesgue se escala a través del determinante de la matriz Jacobiana del sistema, es decir, $dy = \det(C^{1/2}) du = \sqrt{\det(C)} \, du$. Ajustando consecuentemente el dominio y el argumento de la f.d.p. estándar de $Z$, el sumatorio se transforma en:
\begin{align*}
    &= \int_{E} f_Z\left(C^{-1/2}(y-m)\right) \frac{dy}{\sqrt{\det(C)}} = \int_{E} \frac{1}{(2\pi)^{n/2} \sqrt{\det(C)}} \exp\left[ -\frac{1}{2} \left\| C^{-1/2}(y-m) \right\|^2 \right] \, dy = 
\end{align*}
Por último, aplicando las propiedades de la transposición de operadores matriciales y la linealidad del producto escalar, la norma cuadrada del argumento linealizado se reestructura como una forma bilineal gobernada por la matriz de precisión inversa:
\[
\left\| C^{-1/2}(y-m) \right\|^2 = \langle C^{-1/2}(y-m), \, C^{-1/2}(y-m) \rangle = \langle y-m, \, \left(C^{-1/2}\right)^T C^{-1/2}(y-m) \rangle = \langle y-m, \, C^{-1}(y-m) \rangle
\]
Sustituyendo este núcleo en la integral del volumen medible, concluimos la densidad de probabilidad espacial:
\begin{align*}
    &= \int_{E} \frac{1}{(2\pi)^{n/2} \sqrt{\det(C)}} e^{-\frac{1}{2} \langle y-m, \, C^{-1}(y-m) \rangle} \, dy
\end{align*}

\noindent \textbf{Conclusión del Teorema:} \\
Puesto que la relación probabilística se verifica formalmente para cualquier caja elemental $E \subset \mathbb{R}^n$, el integrando resultante se identifica de manera unívoca como el núcleo de la densidad conjunta del sistema. Así, queda demostrado con total rigor que la f.d.p. del vector Gaussiano generalizado $\underline{X}$ es:
\[
f_{\underline{X}}(x) = \frac{1}{(2\pi)^{n/2} \sqrt{\det(C)}} \exp\left[ -\frac{1}{2} \langle x-m, \, C^{-1}(x-m) \rangle \right]
\]

\begin{ejemplo}[Deducción Explícita para la Dimensión $n = 2$]
Desarrollemos rigurosamente la función de densidad de probabilidad (f.d.p.) de un vector Gaussiano multivariante en el caso bidimensional. Fijemos el orden del sistema en $n = 2$, definiendo el vector aleatorio de componentes continuas $\underline{X}$ y su correspondiente vector de evaluación espacial $x$ como:
\[
{\color{blue}\underline{X} = (X_1, X_2) \in \mathbb{R}^2, \qquad x = (x_1, x_2) \in \mathbb{R}^2}
\]
Supongamos que el vector de medias o valores esperados del sistema viene dado por $m = (m_1, m_2) = (\mathbb{E}[X_1], \mathbb{E}[X_2])$.

\vspace{0.4cm}

\noindent \textbf{1. Parametrización y Determinante de la Matriz de Covarianza} \\
La matriz de covarianza simétrica y definida positiva $C$ se construye a partir de las varianzas marginales y las covarianzas cruzadas de las componentes del vector:
\[
C = \begin{pmatrix} \text{Var}(X_1) & \text{Cov}(X_1, X_2) \\ \text{Cov}(X_2, X_1) & \text{Var}(X_2) \end{pmatrix} = \begin{pmatrix} \sigma_{X_1}^2 & \sigma_{X_1 X_2} \\ \sigma_{X_1 X_2} & \sigma_{X_2}^2 \end{pmatrix} = \begin{pmatrix} \sigma_1^2 & \sigma_1 \sigma_2 \rho_{12} \\ \sigma_1 \sigma_2 \rho_{12} & \sigma_2^2 \end{pmatrix}
\]
Donde introducemos las simplificaciones estadísticas estandarizadas para las desviaciones típicas ($\sigma_{X_1} = \sigma_1$, $\sigma_{X_2} = \sigma_2$), la covarianza ($\sigma_{X_1 X_2} = \sigma_{12}$) y el coeficiente de correlación lineal de Pearson $\rho_{12}$, definido axiomáticamente como:
\[
\rho_{12} = \rho_{X_1 X_2} = \frac{\sigma_{X_1 X_2}}{\sigma_{X_1} \sigma_{X_2}} = \frac{\sigma_{12}}{\sigma_1 \sigma_2}
\]
Calculemos el determinante de la matriz de covarianza $C$ para estructurar el factor de normalización de la densidad:
\[
\det(C) = \det \begin{pmatrix} \sigma_1^2 & \sigma_1 \sigma_2 \rho_{12} \\ \sigma_1 \sigma_2 \rho_{12} & \sigma_2^2 \end{pmatrix} = \sigma_1^2 \sigma_2^2 - (\sigma_1 \sigma_2 \rho_{12})^2 = \sigma_1^2 \sigma_2^2 (1 - \rho_{12}^2)
\]
Extrayendo la raíz cuadrada del determinante, obtenemos el coeficiente del denominador de la constante de normalización:
\[
\sqrt{\det(C)} = \sigma_1 \sigma_2 \sqrt{1 - \rho_{12}^2}
\]

\noindent \textbf{2. Inversión Algebraica de la Matriz de Covarianza ($C^{-1}$)} \\
Para hallar la matriz inversa de precisión o concentración, planteamos la identidad fundamental del producto matricial $C \cdot C^{-1} = I_2$:
\[
\begin{pmatrix} \sigma_1^2 & \sigma_1 \sigma_2 \rho_{12} \\ \sigma_1 \sigma_2 \rho_{12} & \sigma_2^2 \end{pmatrix} \begin{pmatrix} a & b \\ c & d \end{pmatrix} = \begin{pmatrix} 1 & 0 \\ 0 & 1 \end{pmatrix} 
\]
Esta multiplicación matricial da origen a un sistema lineal acoplado de cuatro ecuaciones algebraicas con cuatro incógnitas ($a, b, c, d$):
\[
\begin{cases}
    1) \ \sigma_1^2 a + \sigma_1 \sigma_2 \rho_{12} c = 1 \\
    2) \ \sigma_1^2 b + \sigma_1 \sigma_2 \rho_{12} d = 0 \\
    3) \ \sigma_1 \sigma_2 \rho_{12} a + \sigma_2^2 c = 0 \\
    4) \ \sigma_1 \sigma_2 \rho_{12} b + \sigma_2^2 d = 1
\end{cases}
\]

\noindent $\bullet$ \textbf{Resolución de la segunda columna ($b, d$):} \\
Aislando el coeficiente $b$ a partir de la relación lineal de la ecuación 2):
\[
b = -\frac{\sigma_1 \sigma_2 \rho_{12}}{\sigma_1^2} d = -\frac{\sigma_2}{\sigma_1} \rho_{12} d
\]
Sustituyendo directamente esta equivalencia en la ecuación 4) para romper el acoplamiento:
\[
\sigma_1 \sigma_2 \rho_{12} \left( -\frac{\sigma_2}{\sigma_1} \rho_{12} d \right) + \sigma_2^2 d = 1 \implies -\sigma_2^2 \rho_{12}^2 d + \sigma_2^2 d = 1 \implies \sigma_2^2 (1 - \rho_{12}^2) d = 1
\]
\[
\implies \framebox{$\displaystyle d = \frac{1}{\sigma_2^2 (1 - \rho_{12}^2)}$}
\]
Sustituyendo de vuelta el valor obtenido para $d$ en la expresión despejada de $b$, deducimos:
\[
b = -\frac{\sigma_2}{\sigma_1} \rho_{12} \left( \frac{1}{\sigma_2^2 (1 - \rho_{12}^2)} \right) \implies \framebox{$\displaystyle b = -\frac{\rho_{12}}{\sigma_1 \sigma_2 (1 - \rho_{12}^2)}$}
\]

\vspace{0.3cm}

\noindent $\bullet$ \textbf{Resolución de la primera columna ($a, c$):} \\
Aislando de forma homóloga el coeficiente $c$ a partir de la relación lineal de la ecuación 3):
\[
c = -\frac{\sigma_1 \sigma_2 \rho_{12}}{\sigma_2^2} a = -\frac{\sigma_1}{\sigma_2} \rho_{12} a
\]
Sustituyendo de manera análoga esta expresión en la ecuación 1):
\[
\sigma_1^2 a + \sigma_1 \sigma_2 \rho_{12} \left( -\frac{\sigma_1}{\sigma_2} \rho_{12} a \right) = 1 \implies \sigma_1^2 a - \sigma_1^2 \rho_{12}^2 a = 1 \implies \sigma_1^2 (1 - \rho_{12}^2) a = 1
\]
\[
\implies \framebox{$\displaystyle a = \frac{1}{\sigma_1^2 (1 - \rho_{12}^2)}$}
\]
Sustituyendo el valor obtenido para $a$ en la expresión despejada de $c$, deducimos (lo que confirma la simetría $b=c$):
\[
c = -\frac{\sigma_1}{\sigma_2} \rho_{12} \left( \frac{1}{\sigma_1^2 (1 - \rho_{12}^2)} \right) \implies \framebox{$\displaystyle c = -\frac{\rho_{12}}{\sigma_1 \sigma_2 (1 - \rho_{12}^2)}$}
\]

\vspace{0.4cm}

\noindent Agrupando los cuatro componentes elementales en la estructura matricial original, la inversa de la matriz de covarianza $C^{-1}$ queda determinada de forma exacta como:
\[
C^{-1} = \begin{pmatrix} \dfrac{1}{\sigma_1^2 (1 - \rho_{12}^2)} & -\dfrac{\rho_{12}}{\sigma_1 \sigma_2 (1 - \rho_{12}^2)} \\ -\dfrac{\rho_{12}}{\sigma_1 \sigma_2 (1 - \rho_{12}^2)} & \dfrac{1}{\sigma_2^2 (1 - \rho_{12}^2)} \end{pmatrix}
\]

\strut
\hfill \rule{\textwidth}{0.5pt} \hfill
\strut

\noindent \textbf{3. Desarrollo de la Forma Cuadrática de la Distancia} \\
Planteemos la acción de la matriz de precisión $C^{-1}$ operando sobre el vector de desviaciones centrado $(x-m)$:
\[
C^{-1}(x-m) = \begin{pmatrix} \dfrac{1}{\sigma_1^2 (1 - \rho_{12}^2)} & -\dfrac{\rho_{12}}{\sigma_1 \sigma_2 (1 - \rho_{12}^2)} \\ -\dfrac{\rho_{12}}{\sigma_1 \sigma_2 (1 - \rho_{12}^2)} & \dfrac{1}{\sigma_2^2 (1 - \rho_{12}^2)} \end{pmatrix} \begin{pmatrix} x_1 - m_1 \\ x_2 - m_2 \end{pmatrix}
\]
\[
= \begin{pmatrix} \dfrac{1}{\sigma_1^2 (1 - \rho_{12}^2)}(x_1 - m_1) - \dfrac{\rho_{12}}{\sigma_1 \sigma_2 (1 - \rho_{12}^2)}(x_2 - m_2) \\[0.4cm] \dfrac{1}{\sigma_2^2 (1 - \rho_{12}^2)}(x_2 - m_2) - \dfrac{\rho_{12}}{\sigma_1 \sigma_2 (1 - \rho_{12}^2)}(x_1 - m_1) \end{pmatrix}
\]
Para construir el argumento exponencial de la densidad multivariante, ejecutamos el producto escalar interno $\langle x-m, \, C^{-1}(x-m) \rangle$, premultiplicando el vector columna anterior por el vector fila de desviaciones $(x_1 - m_1, \, x_2 - m_2)$:
\begin{align*}
&\langle x-m, \, C^{-1}(x-m) \rangle = \\
&= (x_1 - m_1, \, x_2 - m_2) \begin{pmatrix} \dfrac{1}{\sigma_1^2 (1 - \rho_{12}^2)}(x_1 - m_1) - \dfrac{\rho_{12}}{\sigma_1 \sigma_2 (1 - \rho_{12}^2)}(x_2 - m_2) \\[0.4cm] \dfrac{1}{\sigma_2^2 (1 - \rho_{12}^2)}(x_2 - m_2) - \dfrac{\rho_{12}}{\sigma_1 \sigma_2 (1 - \rho_{12}^2)}(x_1 - m_1) \end{pmatrix} = 
\end{align*}
Expandiendo el álgebra distributiva del producto y agrupando los dos términos cruzados idénticos en un único factor multiplicativo doble, la forma cuadrática se reduce a la expresión escalar continua:
\begin{align*}
&= \frac{1}{\sigma_1^2 (1 - \rho_{12}^2)} (x_1 - m_1)^2 - 2 \frac{\rho_{12}}{\sigma_1 \sigma_2 (1 - \rho_{12}^2)} (x_1 - m_1)(x_2 - m_2) + \frac{1}{\sigma_2^2 (1 - \rho_{12}^2)} (x_2 - m_2)^2
\end{align*}

\noindent \textbf{4. Formulación Exclícita Final de la Densidad Bivariante} \\
Sustituyendo el término del determinante extraído $\sqrt{\det(C)}$ en la constante multiplicativa externa $\frac{1}{(2\pi)^{n/2}\sqrt{\det(C)}}$ para $n=2$, y el polinomio escalar expandido en la componente interna del exponente $-\frac{1}{2}\langle \cdot \rangle$, la función de densidad de probabilidad conjunta bivariante $f_{X_1, X_2}(x_1, x_2)$ se escribe analíticamente como:
\begin{align*}
f_{X_1, X_2}(x_1, x_2) = \frac{1}{2\pi \sigma_1 \sigma_2 \sqrt{1 - \rho_{12}^2}} \exp \Biggl( &-\frac{(x_1 - m_1)^2}{2\sigma_1^2 (1 - \rho_{12}^2)} - \frac{(x_2 - m_2)^2}{2\sigma_2^2 (1 - \rho_{12}^2)} \ + \\
&+ \frac{\rho_{12}}{\sigma_1 \sigma_2 (1 - \rho_{12}^2)} (x_1 - m_1)(x_2 - m_2) \Biggr)
\end{align*}

\end{ejemplo}

\section{Existencia de Procesos Gaussianos}

\noindent
Hasta ahora hemos trabajado con vectores Gaussianos de dimensión finita. El siguiente paso es preguntarnos si, dadas una función de medias y una función de covarianzas "razonables", existe realmente un \textbf{proceso} Gaussiano (infinitas variables, una por cada $t\ge0$) que las tenga como tales.

\begin{definicion}[Función simétrica y positiva definida]
Una función $K:\mathbb{R}^+\times\mathbb{R}^+\to\mathbb{R}$ es:
\begin{itemize}
    \item \textbf{Simétrica} si $K(t,s)=K(s,t)$ para todo $s,t$.
    \item \textbf{Positiva definida} si para todo $n\in\mathbb{N}$, todo $0\le t_1<\dots<t_n$ y todo $c_1,\dots,c_n\in\mathbb{R}$ se tiene $\sum_{i,j} c_ic_j K(t_i,t_j)\ge 0$.
\end{itemize}
\end{definicion}

\begin{definicion}[Proceso Gaussiano, media y covarianza]
Sea $(X_t)_{t\ge0}$ un proceso Gaussiano en $(\Omega,\mathcal{F},P)$ (es decir, todo vector finito $(X_{t_1},\dots,X_{t_n})$ es Gaussiano). Se definen:
\[
K(s,t) = \text{Cov}(X_s,X_t) \quad \text{(función de covarianza)}, \qquad m(t)=E[X_t] \quad \text{(función de media)}.
\]
El proceso es \textbf{centrado} si $m(t)=0\ \forall t$; en general podemos centrarlo con $Y_t = X_t - m(t)$.
\end{definicion}

\begin{prop}
La función de covarianza $K(s,t)$ de un proceso Gaussiano es siempre simétrica y positiva definida (es consecuencia directa de que $\text{Var}\left(\sum_i c_iX_{t_i}\right)\ge0$).
\end{prop}

\begin{teo}[Existencia, vía Kolmogorov]
Sean $m:\mathbb{R}^+\to\mathbb{R}$ una función arbitraria y $K:\mathbb{R}^+\times\mathbb{R}^+\to\mathbb{R}$ simétrica y positiva definida. Entonces existe un espacio de probabilidad $(\Omega,\mathcal{F},P)$ y un proceso Gaussiano $(X_t)_{t\ge0}$ sobre él tal que $E[X_t]=m(t)$ y $\text{Cov}(X_s,X_t)=K(s,t)$.
\end{teo}

\begin{observacion}[Idea de la demostración]
Basta construir, para cada conjunto finito de tiempos $S=\{t_1,\dots,t_K\}$, la matriz $K^S=(K(t_i,t_j))_{i,j}$, que es simétrica y positiva definida y por tanto define un vector Gaussiano $X^S$ válido. Como estas matrices son compatibles entre sí (coinciden en las entradas comunes cuando $S\subset S'$), se cumplen las \textbf{condiciones de consistencia de Kolmogorov}, y el \textbf{Teorema de Extensión de Kolmogorov} garantiza que existe una única medida de probabilidad sobre el proceso completo $(X_t)_{t\ge0}$ con esas distribuciones finito-dimensionales. Este resultado es, esencialmente, el que prueba la \textbf{existencia} de cualquier proceso Gaussiano (en particular, del Movimiento Browniano).
\end{observacion}

\section{Estacionariedad}

\begin{definicion}[Estacionariedad]
Un proceso real $(X_t)_{t\ge0}$ es:
\begin{itemize}
    \item \textbf{(Estrictamente) estacionario} si para todo $n\in\mathbb{N}$, todo $0\le t_1<\dots<t_n$ y todo $s>0$, los vectores $(X_{t_1},\dots,X_{t_n})$ y $(X_{t_1+s},\dots,X_{t_n+s})$ tienen la misma distribución.
    \item \textbf{Débilmente estacionario} si $E[X_t]=E[X_{t+s}]$ y $\text{Cov}(X_{t_1},X_{t_2}) = \text{Cov}(X_{t_1+s},X_{t_2+s})$ para todo $t_1,t_2,s\ge0$.
\end{itemize}
\end{definicion}

\begin{observacion}
Definiendo la \textbf{función de autocovarianza} $c(t,t+s)=\text{Cov}(X_t,X_{t+s})$, la estacionariedad débil equivale a que $c$ dependa \textbf{solo de la diferencia} $s$ entre los tiempos, no de su valor absoluto: $c(u,u+v)=c(0,v)$.
\end{observacion}

\begin{teo}
Si $(X_t)_{t\ge0}$ es un proceso \textbf{Gaussiano}, entonces es estrictamente estacionario si y solo si es débilmente estacionario.
\end{teo}
\begin{observacion}
Esta equivalencia es \textbf{especial de los procesos Gaussianos}: fuera del contexto Gaussiano, en general la estacionariedad débil no implica la fuerte (la débil solo controla los dos primeros momentos, mientras que para procesos Gaussianos la distribución completa queda determinada por la media y la covarianza).
\end{observacion}

\section{Propiedades de las Trayectorias}

\begin{definicion}[Indistinguibilidad y versiones]
Sean $(X_t)_{t\ge0}$ e $(Y_t)_{t\ge0}$ procesos reales en $(\Omega,\mathcal{F},P)$.
\begin{itemize}
    \item Son \textbf{indistinguibles} si existe un conjunto de medida nula $N$ tal que $X_t(\omega)=Y_t(\omega)$ para todo $\omega\notin N$ y \textbf{todo} $t\ge0$ simultáneamente (condición muy fuerte: un único $N$ vale para todos los tiempos a la vez).
    \item Son \textbf{versiones} (o modificaciones) la una de la otra si para cada $t\ge0$ existe un conjunto de medida nula $N_t$ (que puede depender de $t$) tal que $X_t(\omega)=Y_t(\omega)$ para todo $\omega\notin N_t$.
\end{itemize}
\end{definicion}

\begin{definicion}[Continuidad de Hölder]
Una función $\rho:\mathbb{R}^+\to\mathbb{R}$ es \textbf{localmente Hölder-continua} en $t\ge0$ con exponente $\gamma>0$ si existen $\varepsilon,c>0$ tales que $|\rho(t)-\rho(s)|\le c|t-s|^\gamma$ para todo $s$ con $|t-s|<\varepsilon$.
\end{definicion}

\begin{teo}[Kolmogorov-Chentsov]
Sea $(X_t)_{t\ge0}$ un proceso real en $(\Omega,\mathcal{F},P)$. Supongamos que existen $\alpha\ge1,\ \beta,c>0$ tales que
\[
E\big[|X_t-X_s|^\alpha\big] \le c\,|t-s|^{1+\beta} \qquad \forall\, s,t\ge0.
\]
Entonces, para cada $T>0$, $(X_t)$ tiene una versión $(Y_t)_{t\in[0,T]}$ con trayectorias $\gamma$-Hölder continuas, para cualquier $0\le \gamma < \beta/\alpha$.
\end{teo}

\begin{teo}[Corolario para procesos Gaussianos centrados]
Sea $(X_t)_{t\ge0}$ un proceso Gaussiano centrado con función de covarianza $K(s,t)$. Si existen $\beta,c>0$ tales que
\[
K(t,t)+K(s,s)-2K(t,s) \le c\,|t-s|^\beta \qquad \forall\, s,t\ge0,
\]
entonces para todo $\gamma<\beta/2$ el proceso tiene una versión con trayectorias $\gamma$-Hölder continuas.
\end{teo}
\begin{observacion}[Idea de la demostración]
El término de la izquierda es exactamente $E[(X_t-X_s)^2]=\text{Var}(X_t-X_s)$. Usando que los momentos pares de una Gaussiana centrada se expresan en función de la varianza, $E[(X_t-X_s)^{2n}] = \frac{(2n)!}{n!2^n}\big[\text{Var}(X_t-X_s)\big]^n \le c'|t-s|^{(n\beta-1)+1}$, y aplicando Kolmogorov-Chentsov con $\alpha=2n$ para cada $n$, se obtiene Hölder continuidad para exponente $\gamma<\dfrac{n\beta-1}{2n}$; optimizando en $n\to\infty$ se llega a $\gamma<\beta/2$.
\end{observacion}

\begin{definicion}[Recordatorio: Propiedad de Markov]
Un proceso $(X_t)_{t\ge0}$ tiene la \textbf{propiedad de Markov} si para todo $n\in\mathbb{N}$, todo $t_1<\dots<t_n$ y todo $x,x_1,\dots,x_{n-1}\in\mathbb{R}$,
\[
P(X_{t_n}\le x \mid X_{t_1}=x_1,\dots,X_{t_{n-1}}=x_{n-1}) = P(X_{t_n}\le x \mid X_{t_{n-1}}=x_{n-1}).
\]
\end{definicion}

\section{RKHS y Expansión de Karhunen-Loève}

\noindent
Estos dos resultados se enuncian a nivel de cultura general (no se exige reproducir las demostraciones, que son largas): permiten representar cualquier proceso Gaussiano centrado como una "serie de Fourier" aleatoria.

\begin{teo}[Reproducing Kernel Hilbert Space, RKHS]
Sea $K:\mathbb{R}^+\times\mathbb{R}^+\to\mathbb{R}$ una función de covarianza continua. Existe un espacio de Hilbert $\mathcal{H}_K$ de funciones reales continuas, con producto escalar $(\cdot,\cdot)_{\mathcal{H}_K}$, tal que
\[
K(s,\cdot)\in\mathcal{H}_K\ \ \forall s\ge0, \qquad \big(g,K(s,\cdot)\big)_{\mathcal{H}_K} = g(s) \quad \forall\, g\in\mathcal{H}_K,\ s\ge0
\]
(la llamada \textbf{propiedad reproductora}: evaluar $g$ en $s$ equivale a proyectarla escalarmente sobre $K(s,\cdot)$).
\end{teo}

\begin{teo}[Expansión de Karhunen-Loève]
Sea $(X_t)_{t\ge0}$ un proceso Gaussiano centrado con covarianza continua $K$, definido en $(\Omega,\mathcal{F},P)$. Entonces admite la representación, en el sentido de $L^2(\Omega,P)$,
\[
X_t = \sum_{j=1}^{\infty} \varphi_j(t)\, Z_j, \qquad \text{donde } \varphi_j \text{ son funciones continuas en } \mathbb{R}^+ \text{ y } Z_j\overset{i.i.d.}{\sim} N(0,1).
\]
\end{teo}
\begin{observacion}
La idea es construir un isomorfismo entre $\mathcal{H}_K$ y el subespacio de $L^2(\Omega,P)$ generado por las $X_t$; al expresar $K(s,\cdot)$ en una base ortonormal $(\varphi_n)$ de $\mathcal{H}_K$, las variables $Z_j$ resultantes son combinaciones lineales de Gaussianas no correlacionadas, y por tanto \textbf{independientes} entre sí.
\end{observacion}

\section{Movimiento Browniano}

\subsection{Motivación: límite difusivo de un CTRW}

\noindent
Recordamos el modelo de paseo aleatorio en tiempo continuo con $J_i\sim\text{Exp}(\lambda)$ y $X_i\sim N(0,\delta^2)$. Calculando la f.d.p. de $Y(t)$ y pasando al espacio de Laplace-Fourier (vía la fórmula tipo Montroll-Weiss), se obtiene la aproximación
\[
\widehat{p}(u,s) \approx \frac{1}{s+u^2\big(\tfrac{1}{\lambda\delta^2}\big)^{-1}}\cdots
\]
y, tomando el \textbf{límite de escala} $\delta\to0,\ \lambda\to\infty$ manteniendo $\delta^2\lambda=1$ (el "movimiento Browniano emerge cuando los saltos son infinitesimales pero infinitamente frecuentes"), se llega exactamente al núcleo de la ecuación de difusión:
\[
\widehat p(u,s) \longrightarrow \widehat u(u,s) = \frac{1}{s+u^2} \quad\Longrightarrow\quad u(x,t) = \frac{1}{\sqrt{2\pi t}}\, e^{-\frac{x^2}{2t}}.
\]
Esta densidad Gaussiana es la que caracteriza al Movimiento Browniano en cada instante $t$ fijo.

\subsection{Definición y propiedades fundamentales}

\begin{definicion}[Movimiento Browniano Estándar, MBE]
Un \textbf{Movimiento Browniano estándar} (también llamado proceso de Wiener) es un proceso Gaussiano $(B_t)_{t\ge0}$ con
\[
E[B_t]=0\ \ \forall t\ge0, \qquad \text{Cov}(B_s,B_t) = s\wedge t = \min(s,t).
\]
\end{definicion}

\begin{observacion}
La función $K(s,t)=s\wedge t$ es simétrica y positiva definida, así que el Teorema de Existencia (Kolmogorov) garantiza directamente que el MBE \textbf{existe}.
\end{observacion}

\begin{prop}[Distribución de los incrementos]
Para $0\le s<t$, el incremento $B_t-B_s$ es Gaussiano con
\[
E[B_t-B_s]=0, \qquad \text{Var}(B_t-B_s) = E[B_t^2]-2\,\text{Cov}(B_t,B_s)+E[B_s^2] = t-2s+s = t-s,
\]
es decir, $B_t - B_s \sim N(0,t-s)$: la varianza del incremento depende solo de la \textbf{longitud} del intervalo de tiempo.
\end{prop}

\begin{teo}
$(B_t)_{t\ge0}$ admite una versión con trayectorias \textbf{Hölder-continuas} (consecuencia directa del corolario de Kolmogorov-Chentsov, usando que $K(t,t)+K(s,s)-2K(s,t)=|t-s|$, es decir $\beta=1$, lo que da continuidad Hölder para cualquier exponente $\gamma<1/2$).
\end{teo}

\begin{teo}[Propiedades características del MBE]
El Movimiento Browniano Estándar cumple:
\begin{enumerate}
    \item[(i)] $P(B_0=0)=1$.
    \item[(ii)] $t\mapsto B_t(\omega)$ es continua para $P$-casi todo $\omega$.
    \item[(iii)] Para toda partición $0=t_0<t_1<\dots<t_n$, los incrementos $(B_{t_i}-B_{t_{i-1}})_{1\le i\le n}$ son \textbf{independientes} y $B_{t_i}-B_{t_{i-1}}\sim N(0,t_i-t_{i-1})$ (\textbf{incrementos independientes y estacionarios}).
\end{enumerate}
\end{teo}
\begin{observacion}[Idea de la demostración]
(i) y (ii) son consecuencia directa de la definición. Para (iii), la distribución de cada incremento ya se obtuvo arriba; la independencia se prueba viendo que, para $r<s<t$, el vector $(B_s-B_r,\,B_t-B_s)$ es Gaussiano con $\text{Cov}(B_s-B_r,\,B_t-B_s)=0$ (cálculo directo usando $\text{Cov}(B_u,B_v)=u\wedge v$), y por la Proposición de independencia-covarianza-nula vista antes, covarianza cero entre Gaussianas conjuntas implica independencia.
\end{observacion}

\begin{prop}[Propiedades de simetría e invarianza del MBE]
Sea $(B_t)_{t\ge0}$ un MBE. Entonces:
\begin{enumerate}
    \item[(i)] (Traslación espacial) $(x+B_t)_{t\ge0}$ es un Movimiento Browniano que parte de $x$.
    \item[(ii)] (Reflexión) $(-B_t)_{t\ge0} \sim (B_t)_{t\ge0}$.
    \item[(iii)] (Escalado Browniano) Para todo $c>0$, $(\sqrt{c}\,B_{t/c})_{t\ge0} \sim (B_t)_{t\ge0}$.
    \item[(iv)] (Inversión temporal) $Z_t := t\,B_{1/t}$ (con $Z_0:=0$) satisface $(Z_t)_{t\ge0}\sim(B_t)_{t\ge0}$.
    \item[(v)] (Reversión temporal en $[0,t]$) $(B_t)_{t} \sim (B_t - B_{t-\cdot})$.
\end{enumerate}
\end{prop}
\begin{observacion}
Todas se demuestran comprobando que el proceso transformado vuelve a ser Gaussiano, centrado, y con la misma función de covarianza $s\wedge t$ (por ejemplo en (iii): $\text{Cov}(\sqrt{c}B_{s/c},\sqrt{c}B_{t/c}) = c\,(s/c \wedge t/c)= s\wedge t$). La propiedad (iv) es la que permite extender resultados de comportamiento en $t\to0$ a resultados en $t\to\infty$, y viceversa.
\end{observacion}

\begin{prop}[Propiedad de Markov del MBE]
El Movimiento Browniano Estándar es un proceso de Markov.
\end{prop}
\begin{observacion}[Idea de la demostración]
Se define $R(s,t)=\dfrac{K(s,t)}{K(s,s)} = \dfrac{s\wedge t}{s}$, y se comprueba que para $0\le r<s<t$ se cumple la \textbf{factorización} $R(r,t)=R(r,s)R(s,t)$ (de hecho los tres valen $1$ en este caso). Esta condición algebraica es exactamente el \textbf{criterio de Doob} para procesos Gaussianos, que garantiza la propiedad de Markov.
\end{observacion}

\section{Propiedad de Martingala}

\begin{definicion}[Filtración natural]
Para un proceso $(X_t)_{t\ge0}$, la \textbf{filtración natural} $\mathcal{F}_t^X = \sigma(X_s,\, 0\le s\le t)$ es la familia (no decreciente) de $\sigma$-álgebras generadas por la historia del proceso hasta el instante $t$.
\end{definicion}

\begin{definicion}[Martingala]
Sea $(X_t)$ un proceso sobre el espacio filtrado $(\Omega,\mathcal{F},\mathcal{F}_t^X,P)$. Un proceso $(Y_t)$ es una \textbf{martingala} respecto a $(\mathcal{F}_t^X)$ si:
\begin{enumerate}
    \item[(i)] Es \textbf{adaptado} ($Y_t$ es $\mathcal{F}_t^X$-medible) e \textbf{integrable}: $E[|Y_t|]<\infty\ \forall t\ge0$.
    \item[(ii)] $E[Y_t \mid \mathcal{F}_s^X] = Y_s$ para todo $s<t$ (la mejor predicción del futuro, dado el presente, es el valor actual).
\end{enumerate}
\end{definicion}

\begin{lema}[Ya demostrado]
Sea $(B_t)$ un MBE con filtración natural $\mathcal{F}_t^B$. Para $0\le s\le t$, el incremento $B_t-B_s$ es independiente de $\mathcal{F}_s^B$, y por tanto $E[B_t-B_s\mid\mathcal{F}_s^B]=E[B_t-B_s]$.
\end{lema}

\begin{teo}
Sea $(B_t)_{t\ge0}$ un MBE. Entonces tanto $(B_t)$ como $(B_t^2-t)$ son $\mathcal{F}_t^B$-martingalas.
\end{teo}
\begin{observacion}[Idea de la demostración]
Ambos procesos son adaptados (son funciones medibles de $B_t$) e integrables ($E[|B_t|]\le\sqrt{2t/\pi}<\infty$ y $E[|B_t^2-t|]\le 2t<\infty$). Para la propiedad de martingala de $B_t$: usando el lema anterior,
\[
E[B_t\mid\mathcal{F}_s^B] = E[B_t-B_s\mid\mathcal{F}_s^B] + B_s = E[B_t-B_s] + B_s = 0+B_s = B_s.
\]
Para $B_t^2-t$, escribiendo $B_t=(B_t-B_s)+B_s$ y usando de nuevo la independencia del incremento respecto a $\mathcal{F}_s^B$:
\[
E[B_t^2\mid\mathcal{F}_s^B] = E[(B_t-B_s)^2] + B_s^2 + 2B_s\underbrace{E[B_t-B_s\mid\mathcal{F}_s^B]}_{=0} = (t-s)+B_s^2,
\]
luego $E[B_t^2-t\mid\mathcal{F}_s^B] = B_s^2-s$.
\end{observacion}

\begin{teo}[Caracterización de Lévy del Movimiento Browniano]
Sea $(X_t)_{t\ge0}$ un proceso real con filtración natural $(\mathcal{F}_t^X)$. Si:
\begin{enumerate}
    \item[(i)] $X_0=0$ c.s.,
    \item[(ii)] $t\mapsto X_t$ es continua,
    \item[(iii)] tanto $(X_t)$ como $(X_t^2-t)$ son $\mathcal{F}_t^X$-martingalas,
\end{enumerate}
entonces $(X_t)_{t\ge0}$ es un Movimiento Browniano Estándar.
\end{teo}

\begin{observacion}
Este teorema es muy potente en la práctica: \textbf{es el recíproco} del resultado anterior, y permite \textbf{verificar} que un proceso dado (por ejemplo, construido como solución de una ecuación diferencial estocástica) es un Movimiento Browniano sin tener que comprobar directamente la definición Gaussiana, sino solo estas tres condiciones (más fáciles de chequear en la práctica).
\end{observacion}